\documentclass[11pt]{scrartcl}

\usepackage[stixtwo]{fontsetup}
\usepackage{amsmath}
\usepackage{tikz}

\begin{document}
Az $i$ pont körüli forgatás: eltoljuk a síkot úgy, hogy az $i$ az origóba
kerüljön, az origó körül elforgatunk ($e^{i\varphi}$-vel szorzunk), majd
visszatoljuk. A $45^\circ$-os negatív irányú forgatás az
$e^{-i \frac{\pi}{4}}$-gyel való szorzás, így
\[
	g(z)
	= i + e^{-i \frac{\pi}{4}} (z - i)
	= i + \left( \frac{\sqrt{2}}{2} - \frac{\sqrt{2}}{2} i \right) (z - i).
\]

Például a $z = 0$ pont képe
$g(0) = i - i e^{-i \frac{\pi}{4}}
	= -\frac{\sqrt{2}}{2} + \left(1 - \frac{\sqrt{2}}{2}\right) i$:
\begin{center}
	\begin{tikzpicture}[scale=1.6, >=stealth]
		\draw[->] (-1.6, 0) -- (1.4, 0) node[right] {$\operatorname{Re}$};
		\draw[->] (0, -0.4) -- (0, 2.2) node[above] {$\operatorname{Im}$};
		\draw[dashed] (0, 1) circle (1);
		\draw (0, 1) -- (-0.707, 0.293);
		\draw[->] (0, 0.6) arc[start angle=270, end angle=225, radius=0.4];
		\node at (-0.3, 0.45) {$45^\circ$};
		\fill (0, 1) circle (2pt) node[right] {$i$};
		\fill (0, 0) circle (2pt) node[below right] {$z = 0$};
		\fill (-0.707, 0.293) circle (2pt) node[left] {$g(0)$};
	\end{tikzpicture}
\end{center}
\end{document}
